\subsection{Sheaf theory}



In the following, $X$ is a topological space and $\Tau_X$ will denote its topology category. 

\begin{mydef}

Let $\cat{C}$ be a category (often $\cat{C} = \mathscr{Set}, \mathscr{Ab}, \mathscr{Rings}$...).\\
A \textit{presheaf} of $\cat{C}$ is a contravariant functor $\sheaf : \Tau_X \to \cat{C}$.\\
Explicitly, it is the data of :
\begin{itemize}
\item[1.] for every open set $U$ of $X$, an object $\sheaf(U)$ of $\cat{C}$,
\item[2.] for every couple of open sets $V \subseteq U$ of $X$, a morphism $\res_{U,V} : \sheaf(V) \to \sheaf(U)$ verifying :
\begin{itemize}
\item[a.] for every open set $U$ of $X$, then $\res_{U,U} = \text{id}_U$,
\item[b.] for every open sets $U\subseteq V \subseteq W$ of $X$, then $\res_{U,V} \circ \res_{V,W} = \res_{U,W}$.
\end{itemize}
\end{itemize}
The objects $\sheaf(U)$ are often called \textit{sections of $\sheaf$ over $U$}, and for avoiding confusion, will we denoted $\Gamma(U,\sheaf)$.
\end{mydef}

\begin{myrem}
If the context is clear, we will denote $\text{res}_{U,V}(f) = f|_U$ (the notation is justified in the next lines).
\end{myrem}

\begin{myex}
\begin{itemize}
\item[1)] The sheaf of real valued continuous functions $C^0(\--,\RR)$ with classical restriction functions $\res(U,V)(f) = f|_U$.
\item[2)] If $X$ has an additional structure, the above example can be further specified : if $X$ is a differentiable manifold, we can consider $C^\infty(\--,\RR)$, if $X$ is a complex manifold, $\holo(\--,\CC)$ ...
\item[3)] If $C$ is an object of $\cat{C}$, the constant presheaf $\underline{C}$. We will be intersted in $\underline{\ZZ}$ or $\underline{\CC}$.
\item[4)] The sheaf of real valued bounded functions $B(\--,\RR)$.
\end{itemize}
\end{myex}

Presheaves comes from a category point of view, so we need morphisms :

\begin{mydef}
A \textit{morphism of presheaves} (over the same space $X$) $\varphi : \sheaf \to \mathcal{G}$ is a natural transformation.\\
More explicitly, it is the data, for each $U\in \Tau_X$, of a morphism $\varphi_U : \sheaf(U) \to \mathcal{G}(U)$ of $\cat{C}$ such that the following diagram commutes for every $U \subseteq V$ :
\begin{center}
\begin{tikzcd}
\sheaf(V) \arrow[r, "\varphi_V"] \arrow[d, "\res^{\sheaf}_{U,V}" left]
& \mathcal{G}(V) \ \arrow[d, "\res^{\mathcal{G}}_{U,V}"] \\
\sheaf(U) \arrow[r, "\varphi_U"]
& \mathcal{G}(U)
\end{tikzcd}
\end{center}
\end{mydef}


\begin{mydef}
A \textit{sheaf} $\sheaf$ is a presheaf satisfying the extra conditions for any open cover $\{U_i,i\in I\}$ of an open set $U$ of $X$:
\begin{itemize}
\item[3.] for every $f,g \in \sheaf(U)$ such that : $\forall i \in I,~ f|_{U_i} = g|_{U_i}$, then $f = g$.
\item[4.] for every $(f_i)\in \prod_{i\in I} \sheaf(U_i)$ such that : $\forall (i,j)\in I^2,~ f_i|_{U_i \cap U_j}= f_j|_{U_i \cap U_j}$, then there exists $f\in\sheaf(U)$ with $f|_{U_i} = f_i$ for any $i\in I$.
\end{itemize}
\end{mydef}


\begin{myrem}
For $U = \emptyset$, it says that $\sheaf(\emptyset)$ is the terminal object of $\cat{C}$ (if such one exists). Indeed, we can take the open covering of $\emptyset$ by a family indexed by $\emptyset$ itself, and we then have an arrow $\sheaf(\emptyset) \to $
\end{myrem}

\begin{myex}
\begin{itemize}
\item 1) and 2) are indeed sheaves.
\item 3) is not, because of the previous remark.
\item 4) is not either, with the conter example $\RR = \bigcup_{n\in \NN} ~(-n,n)$ and $f : x\mapsto x$.
\item If $X$ is Hausdorff, $x_*$ is a point of $X$, $(S,s)$ a pointed set, the \textit{skycraper sheaf} is given by $\sky_{x_*}(S) : U \mapsto \left\lbrace
\begin{array}{l l}
S & \text{if } x_* \in U \\
s & \text{otherwise}.
\end{array}
\right.
$
\end{itemize}
\end{myex}

\begin{mydef}
We will denote by $\text{PSh}_X$, respectively $\text{Sh}_X$, the category of presheaves, respectively the full subcategory of sheaves, over $X$.
\end{mydef}

One can ask what looks like the behavior of the sheaf very locally, on a point. Since a singleton is in general not a point, we need the following definition.

\begin{mydef}
Let $\sheaf$ be a (pre)sheaf.\\
We define the \textit{stalk} at $x\in X$ as :
\begin{center}
$\sheaf_x := \quot{\{(U,s) | ~x\in U, ~U \in \Tau_X, ~s \in \sheaf(U)\}}{\sim}~,$
\end{center}
where $(U,s) \sim (V,t) \Leftrightarrow (\exists W \subset U \cap V): s|_W = t|_W$.
\end{mydef}

\begin{myrem}
From a categorical point of view, it is also the direct limit : $\sheaf_x = \underset{\overrightarrow{U \ni x}}{\lim} \sheaf(U)$, where the limit is taken over open set $U$.
\end{myrem}

It turns out that this localization principle can be used to transform a presheaf into a sheaf in a natural way (in the sense of what follows).

\begin{mydef}
Let $\sheaf$ be a presheaf.\\
We define for every open set $U$
\begin{center}
$\d \sheaf^{\sharp}(U) = \left\lbrace s : U \to \coprod_{x\in U} \sheaf_x ~\left|\right.~ \left\lbrace 
\begin{array}{l}
\forall x\in U,~ s(x) \in \sheaf_x\\
\exists V \subset U,\exists t\in \sheaf(V) \text{ s.th.} \forall y\in V,~ s(y) = t(y)
\end{array}
\right. \right\rbrace$.
\end{center}
$\sheaf^\sharp$ is called the \textit{sheafification of $\sheaf$}.
\end{mydef}

\begin{myprop}
$\sheaf^\sharp$ is a sheaf.
\end{myprop}

\begin{myex}
$\underline{\RR}^\sharp$ is the sheaf of locally constant functions, more precisely the sheaf of continuous functions $C^0(\--,\RR)$, where $\RR$ is equipped with the discrete topology.
\end{myex}



\begin{myprop}
$((\--)^\sharp,\text{Forget})$ is an adjoint pair of functors between $\text{PSh}_X$ and $\text{Sh}_X$.
\end{myprop}

\begin{mydemo}
.
\end{mydemo}

\begin{myex}
Suppose that $\cat{C}$ admits kernels and cokernels (this assumption holds whenever kernels and cokernels are written). \\
For $\varphi: \sheaf \to \mathcal{G}$, we define $\Ker(\varphi)(U) := \Ker(\varphi_U)$, $\text{Im}(\varphi)(U) := \text{Im}(\varphi_U)$ and $\text{Coker}(\varphi)(U) := \text{Coker}(\varphi_U)$.\\
The first one is always a sheaf, but the two others are not in general. In the following, $\text{Im}(\varphi)$ or $\text{Coker}(\varphi)$ will always denote the sheafification of these presheaves.
\end{myex}

\begin{myprop}
A complex $0 \to \sheaf \to \mathcal{G} \to \mathcal{H} \to 0$ is exact if, and only if the induced complex on stalks $0 \to \sheaf_x \to \mathcal{G}_x \to \mathcal{H}_x \to 0$ is exact for every $x\in X$.
\end{myprop}

\begin{mydemo}
.
\end{mydemo}

\begin{myrem}
In general, it is not equivalent of the exactness of the induced sequence over every open set $U$.\\
One could take as a counterexample the exponential sequence on $\CC$ :
\begin{center}
$
\begin{array}{c c c c c c c c c}
0 & \to & \underline{\ZZ} & \to & \holo_\CC & \to & \holo_\CC^* & \to & 0,\\
  &     &                &     &   f       & \mapsto & \exp(2\pi i f) & &\\
\end{array}
$
\end{center}
that is not exact over $\CC^*$, because the map $\holo_{\CC^*} \to \holo_{\CC^*}^*$ is not surjective : the map $z \mapsto z$ does not belong to the image since $\log$ is not well defined over $\CC^*$.
\end{myrem}


\begin{myprop} \label{patching_sheaf}
Let $X = \cup_{i\in I} U_i$ an open covering of $X$. For every $i\in I$, let $\sheaf_i$ a sheaf on $U_i$.\\
Assume that for every $i \neq j \in I$, there exists an isomorphism 
\begin{center}
$\theta_{i,j} : \sheaf_i|_{U_i \cap U_j} \overset{\sim}{\to} \sheaf_j|_{U_i \cap U_j}$,
\end{center}
satisfying the condition $\theta_{j,k} \circ \theta _{i,j} = \theta_{i,k}$ on $U_i \cap U_j \cap U_k$ for every triplet of distinct elements of $I$.\\
Then, there exists a sheaf $\sheaf$ on $X$ such that for every $i\in I, \sheaf|_{U_i} = \sheaf_i$. Moreover, this sheaf is unique, up to unique isomorphism.
\end{myprop}

\begin{mydemo}
.
\end{mydemo}